% Chapter 2 -- Setting Up (Mac OS 8 Installation Ch. 2 form; the operations
% chapter of this book, so it opens with the slogan, operator directive
% 2026-10-17).  Sources (Law 1): spec/hardware.json and spec/memory_map.h
% FLAIR_HEAP_MIN (the requirements, as the FLAIR User's Guide Appendix B);
% os/milton/kmain.c (the startup disk: IO.SYS and INITDOS.SYS, FLAIR started
% with no command; the data disk on the IDE primary slave mounted as drive
% A, label INITECH; TRASH-CREATE and DESKTOP-DB-CREATE on a disk that has
% neither); harness/emu/qemu.c (--disk the startup disk, --disk2 the data
% disk: the first and second drives).
% GATES: test-hardware-spec test-flair-desktop test-flair-data-volume test-desktop-db test-boot-system test-pubs-flair-screens
\Gates{test-hardware-spec, test-flair-desktop, test-flair-data-volume,
  test-desktop-db, test-boot-system, test-pubs-flair-screens}
\chapter{Setting Up}
\IsThisGoodForTheCompany

This chapter tells you what your computer needs, how to put the InitechOS
disks in place, and how to start FLAIR for the first time.

\section{Checking what you need}
\index{system requirements}
To use FLAIR you need:
\begin{itemize}
\item An Initech workstation or a compatible personal computer with an
  Intel 80486 microprocessor.
\item At least 5~MB of memory: the first megabyte and 4~MB of extended
  memory above it.
\item A VGA display adapter with the VESA BIOS Extensions, Version~2.0,
  able to show 640 by 480 dots in 256 colours, and a VGA monitor.
\item A keyboard and a mouse.
\item Two disk drives on an AT-compatible (ATA) controller: one for the
  startup disk and one for the data disk.
\end{itemize}
A math coprocessor is not needed for FLAIR; Initech 123 needs one.

\section{Connecting the computer}
\begin{TaskSteps}
\Step Make sure the computer and the display are turned off.
\Step Connect the keyboard and the mouse to the computer.
\Result{FLAIR uses one mouse button, the left one.}
\Step Connect the display to the VGA adapter, and the display and the
computer to the power.
\end{TaskSteps}
The instructions that came with your computer and your display show where
each cable goes.

\section{Making backup copies}
\index{backup copies}
Before you use the InitechOS disks, make copies of them and keep the
originals in a safe place.  Work with the copies.  If a disk is damaged or a
file is deleted by mistake, you can make a new copy from the original.
\draft{the package and a disk-copying command (DISKCOPY) are not built;
see bead initech-aai2}

\section{Putting the disks in place}
\begin{TaskSteps}
\Step Put the startup disk in the first drive.
\Result{The computer starts InitechDOS and FLAIR from it.}
\Step Put the data disk, labelled INITECH, in the second drive.
\Result{FLAIR shows this disk on the desktop as the \texttt{INITECH} icon,
and InitechDOS calls it drive A.}
\end{TaskSteps}
Leave both disks in place while you use FLAIR.

\section{Starting FLAIR for the first time}
\index{starting FLAIR}
\begin{TaskSteps}
\Step Turn on the display, then the computer.
\Result{The computer checks its memory and loads InitechDOS from the
startup disk.  InitechDOS starts FLAIR by itself.  After a few seconds the
FLAIR desktop appears.}
\end{TaskSteps}

\FlairScreen[label=fig:gs-start]{g_desktop}{FLAIR, ready for work}

The first time FLAIR starts with a data disk, it adds two items to the
disk: the folder in which it keeps the Trash, and the file
\file{DESKTOP.DB}, in which it records the arrangement of your desktop.  You
do not see either of them in the disk's window, and you should leave them
alone.

If the desktop does not appear, or looks different from
Figure~\ref{fig:gs-start}, turn to Chapter~5.
